% 编译方式：XeLaTeX，建议连续编译两遍。
% 本文件使用 Windows 自带的宋体、黑体等中文字体，并将所用字形嵌入 PDF。
\documentclass[UTF8,a4paper,11pt,fontset=windows]{ctexart}
\usepackage[margin=21mm,headheight=15pt,headsep=7mm,footskip=10mm]{geometry}
\usepackage{amsmath,amssymb,booktabs,tabularx,array,enumitem}
\usepackage{xcolor,tikz,tcolorbox,fancyhdr,hyperref}
\usetikzlibrary{arrows.meta,decorations.pathmorphing,patterns}
\definecolor{ink}{HTML}{263747}
\definecolor{accent}{HTML}{246579}
\definecolor{soft}{HTML}{F0F5F7}
\definecolor{muted}{HTML}{5C6872}
\hypersetup{colorlinks=true,linkcolor=accent,urlcolor=accent,
  pdftitle={二阶系统基础复习},pdfauthor={},pdfsubject={小车建模、自然频率、阻尼比、标准二阶模型与谐振}}
\pagestyle{fancy}
\fancyhf{}
\fancyhead[L]{\small\color{muted}基础知识储备\enspace /\enspace 自动控制}
\fancyhead[R]{\small\color{muted}二阶系统}
\fancyfoot[C]{\small\color{muted}\thepage}
\renewcommand{\headrulewidth}{0.3pt}
\renewcommand{\footrulewidth}{0pt}
\ctexset{section={format=\Large\bfseries\color{ink},beforeskip=13pt,afterskip=8pt},
  subsection={format=\normalsize\bfseries\color{ink},beforeskip=9pt,afterskip=5pt}}
\setlength{\parindent}{2em}
\setlength{\parskip}{4pt}
\linespread{1.18}
\setlist[itemize]{leftmargin=1.8em,topsep=3pt,itemsep=2pt,parsep=0pt}
\setlist[enumerate]{leftmargin=1.8em,topsep=3pt,itemsep=3pt,parsep=0pt}
\renewcommand{\arraystretch}{1.3}
\newcolumntype{Y}{>{\raggedright\arraybackslash}X}
\newtcolorbox{keybox}{colback=soft,colframe=accent!30,boxrule=0.4pt,arc=1.2mm,
  left=3mm,right=3mm,top=2mm,bottom=2mm,before skip=7pt,after skip=7pt}
\newcommand{\unit}[1]{\,\mathrm{#1}}
\newcommand{\dd}{\mathrm{d}}
\begin{document}

\begin{center}
{\LARGE\bfseries\color{ink}二阶系统基础复习}\\[5pt]
{\small\color{muted}从小车受力到标准模型\quad |\quad 建模、阻尼性质与谐振}
\end{center}
\vspace{-3pt}

\section{为什么小车模型是二阶系统？}

考虑一辆沿水平方向运动的小车，质量为 $m$，通过弹簧和黏性阻尼器连接到固定墙面。以无外力时的平衡位置为原点，位移为 $x(t)$，向右为正，外力为 $F(t)$。

\begin{center}
\begin{tikzpicture}[x=1cm,y=1cm,>=Latex,line width=0.8pt,font=\small]
  \fill[pattern=north east lines,pattern color=muted!60] (-0.2,-0.15) rectangle (0,2.2);
  \draw (0,-0.15)--(0,2.2);
  \draw[decorate,decoration={coil,aspect=0.45,segment length=5pt,amplitude=4pt}]
    (0,1.5)--(3.2,1.5);
  \node[above,fill=white,inner sep=2pt] at (1.6,1.7) {弹簧 $k$};
  \draw (0,0.65)--(1.3,0.65);
  \draw (1.3,0.4)--(1.3,0.9)--(2.05,0.9);
  \draw (1.3,0.4)--(2.05,0.4);
  \draw (1.75,0.45)--(1.75,0.85);
  \draw (1.75,0.65)--(3.2,0.65);
  \node[below] at (1.65,0.35) {阻尼器 $c$};
  \filldraw[fill=soft] (3.2,0.35) rectangle (5.35,1.95);
  \node at (4.27,1.2) {小车 $m$};
  \draw (3.65,0.16) circle (0.18);
  \draw (4.9,0.16) circle (0.18);
  \draw[muted] (2.9,-0.03)--(7.6,-0.03);
  \draw[->,accent,thick] (5.35,1.2)--(7.3,1.2) node[right,text=ink] {$F(t)$};
  \draw[->] (4.25,2.35)--(6.1,2.35) node[right] {$x$ 正方向};
  \node[below,text=muted] at (4.7,-0.32) {水平方向运动；重力与支持力相互平衡};
\end{tikzpicture}
\end{center}
\vspace{-8pt}

\subsection{从牛顿第二定律建立方程}

采用线性弹簧 $F_k=-kx$ 和线性黏性阻尼 $F_c=-c\dot x$，忽略额外摩擦。牛顿第二定律 $\sum F=m\ddot x$ 给出
\begin{equation}
  F(t)-c\dot x(t)-kx(t)=m\ddot x(t).
\end{equation}
整理得到小车的运动微分方程：
\begin{equation}\label{eq:physical}
  \boxed{m\ddot x(t)+c\dot x(t)+kx(t)=F(t).}
\end{equation}

\noindent
\begin{tabularx}{\linewidth}{@{}lYl@{}}
\toprule
符号 & 物理含义 & 单位 \\
\midrule
$m$ & 质量，表征惯性 & $\mathrm{kg}$ \\
$c$ & 黏性阻尼系数，阻尼力与速度成正比 & $\mathrm{N\,s/m}$ \\
$k$ & 弹簧刚度，恢复力与位移成正比 & $\mathrm{N/m}$ \\
$x,\ \dot x,\ \ddot x$ & 位移、速度、加速度 & $\mathrm{m},\ \mathrm{m/s},\ \mathrm{m/s^2}$ \\
$F(t)$ & 施加在小车上的外力 & $\mathrm{N}$ \\
\bottomrule
\end{tabularx}

\subsection{“二阶”指什么？}
方程中最高阶导数是位移的\textbf{二阶导数} $\ddot x$，因此这是二阶系统。给定外力后，确定完整运动过程还需要两个初始条件：$x(0)$ 和 $\dot x(0)$。仅知道小车在哪里，不足以确定之后怎样运动，还必须知道它此刻的速度。

\begin{keybox}
\textbf{物理直觉：}弹簧把小车拉回平衡位置；惯性使它可能冲过平衡位置；正阻尼消耗机械能，使自由振动逐渐衰减。二阶不表示有两个输入或两个输出。
\end{keybox}

\newpage
\section{自然频率：先看没有阻尼时怎样运动}

\subsection{自然频率\texorpdfstring{ $\omega_n$}{ ωn}：无阻尼时的固有振动节奏}
这里的自然频率 $\omega_n$，准确地说是\textbf{无阻尼自然角频率}：系统没有外部激励、没有阻尼时，由自身参数决定的自由振动角频率。

以小车为例，运动方程为 $m\ddot x+c\dot x+kx=F(t)$。为了考察系统自身的振动，令外力 $F(t)=0$，并暂时忽略阻尼 $c=0$，得到
\[
  m\ddot x+kx=0
  \quad\Rightarrow\quad
  \ddot x+\frac{k}{m}x=0.
\]
假设把小车拉到 $x_0$ 后从静止释放，即 $x(0)=x_0$、$\dot x(0)=0$，它的运动是
\[
  x(t)=x_0\cos\left(\sqrt{\frac{k}{m}}\,t\right).
\]
与正弦振动的形式 $x(t)=x_0\cos(\omega t)$ 对比，便得到
\begin{equation}
  \boxed{\omega_n=\sqrt{\frac{k}{m}}.}
\end{equation}
因此，这个定义有明确的物理含义：
\begin{itemize}
  \item \textbf{刚度 $k$ 越大}，恢复力越强，来回振动越快。
  \item \textbf{质量 $m$ 越大}，惯性越大，来回振动越慢。
  \item 在线性模型中，拉开的距离改变振幅，\textbf{不改变自然频率}。
\end{itemize}

注意单位：$\omega_n$ 的单位是 $\mathrm{rad/s}$。如果问“每秒振动多少次”，应使用单位为赫兹（$\mathrm{Hz}$）的自然频率 $f_n$；自然周期记为 $T_n$：
\[
  \boxed{f_n=\frac{\omega_n}{2\pi}},
  \qquad
  \boxed{T_n=\frac{2\pi}{\omega_n}}.
\]
例如，$m=1\unit{kg}$、$k=4\unit{N/m}$ 时，$\omega_n=2\unit{rad/s}$，即约 $0.318\unit{Hz}$，一个周期约为 $3.14\unit{s}$。

\subsection{有阻尼时，如何理解这个参数？}
\textbf{实际存在阻尼时，$\omega_n$ 仍然是系统的参数，但不一定等于实际振荡频率。}即使阻尼足够大、自由响应不再周期性振荡，$\omega_n$ 仍由质量与刚度按上述公式定义。

要进一步回答“加入阻尼后是否还会振荡、振荡多久消失”，就需要研究含阻尼的自由运动方程。由此可以确定振荡的分界，并引出临界阻尼系数和阻尼比。

\begin{keybox}
\textbf{本节结论：}质量与刚度给出固有的时间尺度 $\omega_n=\sqrt{k/m}$。接下来要确定阻尼相对于这个系统究竟有多强，才能判断完整的响应形态。
\end{keybox}

\newpage
\section{从振荡分界引出临界阻尼与阻尼比}

\subsection{特征根决定自由响应的形式}
恢复阻尼、仍令外力为零，小车满足 $m\ddot x+c\dot x+kx=0$。令 $x=e^{st}$ 代入，可得特征方程及其根（概念与推导见\hyperref[app:characteristic]{附录 A}）：
\begin{equation}\label{eq:characteristic}
  ms^2+cs+k=0,
  \qquad
  s_{1,2}=\frac{-c\pm\sqrt{c^2-4mk}}{2m}.
\end{equation}
这里的 $s$ 是指数响应中的参数。\textbf{实部为负，对应衰减；实部为正，对应增长；非零虚部对应振荡项。}例如，共轭复根 $s=\alpha\pm\mathrm{j}\beta$ 对应的实数响应为
\[
  x_h(t)=e^{\alpha t}\left[C_1\cos(\beta t)+C_2\sin(\beta t)\right],
\]
其中 $\mathrm{j}^2=-1$，$C_1,C_2$ 由初始条件决定。因此判断是否振荡，关键是判别式 $\Delta=c^2-4mk$。

\noindent
\begin{tabularx}{\linewidth}{@{}p{32mm}YY@{}}
\toprule
条件（$c>0$） & 特征根 & 自由响应 \\
\midrule
$\Delta<0$ & 实部为负的一对共轭复根 & 有衰减振荡 \\
$\Delta=0$ & 两个相同的负实根 & 恰好到达非振荡边界 \\
$\Delta>0$ & 两个不同的负实根 & 无周期性振荡 \\
\bottomrule
\end{tabularx}

\subsection{临界阻尼系数来自判别式为零}
\textbf{临界阻尼系数 $c_{\mathrm{cr}}$，是在质量和刚度固定时，使自由响应不再具有振荡项的最小正阻尼系数。}在这一分界处，判别式恰好为零：
\[
  c_{\mathrm{cr}}^2-4mk=0.
\]
取正值，得到
\begin{equation}\label{eq:critical}
  \boxed{c_{\mathrm{cr}}=2\sqrt{mk}=2m\omega_n.}
\end{equation}
它的单位与实际阻尼系数 $c$ 相同，都是 $\mathrm{N\,s/m}$。这个分界同时取决于质量与刚度，因此仅比较 $c$ 的大小，不能判断不同系统属于哪种阻尼状态。

\subsection{阻尼比：实际阻尼是临界阻尼的多少倍}
用临界阻尼作为比较基准，定义无量纲参数
\begin{equation}\label{eq:zeta}
  \boxed{\zeta=\frac{c}{c_{\mathrm{cr}}}=\frac{c}{2\sqrt{mk}}=\frac{c}{2m\omega_n}.}
\end{equation}
于是 $0<\zeta<1$ 为欠阻尼，$\zeta=1$ 为临界阻尼，$\zeta>1$ 为过阻尼。将参数定义反过来写，可以得到后面整理标准模型所需的两个关系：
\begin{equation}\label{eq:coefficients}
  \boxed{\frac{k}{m}=\omega_n^2},
  \qquad
  \boxed{\frac{c}{m}=2\zeta\omega_n}.
\end{equation}

\newpage
\section{不同阻尼比下的系统性质}

\subsection{把特征根写成自然频率和阻尼比的形式}
把式\eqref{eq:coefficients}代入特征方程，得到
\[
  s^2+2\zeta\omega_n s+\omega_n^2=0,
  \qquad
  s_{1,2}=-\zeta\omega_n\pm\omega_n\sqrt{\zeta^2-1}.
\]
下表的\textbf{自由响应}指撤去外力后的运动；\textbf{阶跃响应}指小车从静止开始，外力从 $0$ 突然变为正常数 $F_0$ 后的运动，目标平衡位置为 $F_0/k$。超调表示位移超过这一位置。

\noindent
\begin{tabularx}{\linewidth}{@{}p{24mm}p{17mm}YY@{}}
\toprule
阻尼比 & 名称 & 特征根与自由响应 & 阶跃响应性质 \\
\midrule
$\zeta=0$ & 无阻尼 & 纯虚根；持续等幅振荡 & 持续振荡，不能稳定到目标 \\
\addlinespace[3pt]
$0<\zeta<1$ & 欠阻尼 & 负实部共轭复根；振荡衰减 & 有超调，经过衰减振荡趋于目标 \\
\addlinespace[3pt]
$\zeta=1$ & 临界阻尼 & 重复负实根；无周期性振荡 & 不超调；固定 $\omega_n$ 时，在 $\zeta\geq1$ 的这类模型中趋近目标最快 \\
\addlinespace[3pt]
$\zeta>1$ & 过阻尼 & 两个不同的负实根；两个衰减指数项 & 不超调；固定 $\omega_n$ 时，增大阻尼使响应更迟缓 \\
\addlinespace[3pt]
$\zeta<0$ & 负阻尼 & 根的实部为正；一般扰动增长 & 不稳定，不能收敛到目标 \\
\bottomrule
\end{tabularx}

\subsection{振荡频率与衰减速度分别由什么决定？}
欠阻尼时，自由响应为
\begin{equation}\label{eq:underdamped}
  x_h(t)=e^{-\zeta\omega_n t}\left[C_1\cos(\omega_dt)+C_2\sin(\omega_dt)\right],
  \qquad \omega_d=\omega_n\sqrt{1-\zeta^2}.
\end{equation}
$\omega_d$ 是实际的\textbf{阻尼振荡角频率}，决定振荡节奏；指数包络 $e^{-\zeta\omega_n t}$ 决定振幅衰减速度，包络时间常数为 $1/(\zeta\omega_n)$。

达到临界阻尼时，响应变成
\begin{equation}\label{eq:critical-response}
  x_h(t)=(C_1+C_2t)e^{-\omega_n t}.
\end{equation}
此时已没有正弦、余弦项，但仍需要时间才能回到平衡位置。过阻尼时则是两个衰减指数项的叠加：$x_h(t)=C_1e^{s_1t}+C_2e^{s_2t}$。

\subsection{将分界、分类与物理参数联系起来}
取 $m=1\unit{kg}$、$k=4\unit{N/m}$，则 $\omega_n=2\unit{rad/s}$、$c_{\mathrm{cr}}=4\unit{N\,s/m}$。当 $c$ 分别为 $2$、$4$、$8\unit{N\,s/m}$ 时，$\zeta$ 分别为 $0.5$、$1$、$2$，对应欠阻尼、临界阻尼和过阻尼。

\begin{keybox}
\textbf{适用边界：}$\zeta>0$ 时均渐近稳定；$\zeta=0$ 时振动不衰减，谐振输入还可使输出无界。普通被动阻尼器满足 $c\geq0$，负阻尼表示等效能量注入。“不振荡”指自由响应；持续正弦外力仍可产生受迫振动。
\end{keybox}

\newpage
\section{从参数含义自然得到标准二阶模型}

\subsection{先用动态参数重新表达物理方程}
经过前面的推导，$\omega_n$ 表示固有时间尺度，$\zeta$ 表示相对于临界值的阻尼强弱。现在恢复外力，将式\eqref{eq:physical}除以 $m$：
\[
  \ddot x+\frac{c}{m}\dot x+\frac{k}{m}x=\frac{F(t)}{m}.
\]
代入已经建立的关系 $c/m=2\zeta\omega_n$、$k/m=\omega_n^2$，得到
\begin{equation}\label{eq:force}
  \boxed{\ddot x+2\zeta\omega_n\dot x+\omega_n^2x=\frac{F(t)}{m}.}
\end{equation}
这一步仅仅更换了描述参数，系统与输入仍是原来的小车和外力。标准形式左侧的每个系数，都有前面推导出的物理含义。

\subsection{再统一输入尺度，让稳态输出等于输入}
对 $c>0$ 的稳定小车，恒定外力对应的最终平衡位置为 $x_\infty=F/k$。为了直接用这个平衡位置作为输入，定义
\[
  u(t)=\frac{F(t)}{k},\qquad
  \frac{F(t)}{m}=\frac{k}{m}\frac{F(t)}{k}=\omega_n^2u(t).
\]
将输出位移记为 $y=x$，得到\textbf{稳态增益为 $1$ 的标准二阶模型}：
\begin{equation}\label{eq:standard}
  \boxed{\ddot y+2\zeta\omega_n\dot y+\omega_n^2y=\omega_n^2u.}
\end{equation}
这里的 $u$ 是位移量，表示当前外力所对应的平衡位置；它不表示小车的实际位移。对恒定输入，稳态时导数为零，从而 $y_\infty=u$。右端的 $\omega_n^2$ 来自输入尺度的选择，\textbf{不规定输入必须以频率 $\omega_n$ 变化}。

\subsection{保留稳态增益，得到一般的标准形式}
如果希望保留实际输入输出之间的比例，记稳态增益为 $K$，即恒定输入下 $y_\infty=Ku$，则标准模型写成
\begin{equation}\label{eq:general}
  \boxed{\ddot y+2\zeta\omega_n\dot y+\omega_n^2y=K\omega_n^2u.}
\end{equation}
\noindent
\begin{tabularx}{\linewidth}{@{}YYl@{}}
\toprule
对同一辆小车的输入选择 & 输出与稳态关系 & 稳态增益 \\
\midrule
$u=F$，直接输入外力 & $y=x$，$x_\infty=F/k$ & $K=1/k$ \\
$u=F/k$，输入平衡位置 & $y=x$，$x_\infty=u$ & $K=1$ \\
\bottomrule
\end{tabularx}
\par
大写 $K$ 是稳态增益，小写 $k$ 是弹簧刚度。$K$ 的单位取决于输入输出的选择；外力到位移的增益单位为 $\mathrm{m/N}$。在\textbf{零初始条件}下对式\eqref{eq:general}作拉普拉斯变换，可写成等价的传递函数：
\begin{equation}\label{eq:transfer}
  \boxed{G(s)=\frac{Y(s)}{U(s)}=\frac{K\omega_n^2}{s^2+2\zeta\omega_ns+\omega_n^2}.}
\end{equation}

\begin{keybox}
\textbf{代入检查：}$m=1$、$c=2$、$k=4$（均取上述 SI 单位）时，$\omega_n=2$、$\zeta=0.5$。原方程为 $\ddot x+2\dot x+4x=F$；令 $u=F/4$ 后，标准方程为 $\ddot y+2\dot y+4y=4u$。
\end{keybox}

\newpage
\section{谐振与频率响应}

\subsection{从自由运动转向持续的正弦输入}
前面通过特征根研究了撤去外力后的自由响应。现在改变问题：\textbf{持续施加正弦输入时，哪个输入频率最容易使系统产生大幅响应？}本章采用式\eqref{eq:standard}的单位稳态增益模型，即 $K=1$：
\[
  \ddot y+2\zeta\omega_n\dot y+\omega_n^2y=\omega_n^2u,
  \qquad u(t)=A\sin(\omega t),\quad A>0.
\]
这里 $A$ 是输入幅值，$\omega$ 是可调的\textbf{外部输入角频率}，与系统参数 $\omega_n$ 不同。若 $\zeta>0$，过渡过程会逐渐衰减，最后留下正弦稳态响应：
\begin{equation}\label{eq:sine-steady}
  y_{\mathrm{ss}}(t)=B(\omega)\sin\!\left[\omega t+\varphi(\omega)\right].
\end{equation}
\textbf{稳态输出与输入同频}，改变的是幅值 $B$ 和相位 $\varphi$。比较不同频率下的响应时，应保持输入幅值 $A$ 不变，并比较过渡过程消失后的输出幅值。

\subsection{幅频响应：每个输入频率被放大多少倍？}
在传递函数中令 $s=\mathrm{j}\omega$，并取 $K=1$，得到频率响应
\[
  G(\mathrm{j}\omega)
  =\frac{\omega_n^2}{\omega_n^2-\omega^2+\mathrm{j}\,2\zeta\omega_n\omega}.
\]
复数的模给出幅值增益，辐角给出相位变化。因此
\begin{equation}\label{eq:magnitude}
  \boxed{M(\omega)=\frac{B(\omega)}{A}=|G(\mathrm{j}\omega)|
  =\frac{1}{\sqrt{\left[1-(\omega/\omega_n)^2\right]^2
  +\left(2\zeta\omega/\omega_n\right)^2}}.}
\end{equation}
幅频曲线以\textbf{输入角频率 $\omega$ 为横坐标、幅值增益 $M$ 为纵坐标}。低频极限 $M(0)=1$ 对应静态增益；高频时位移响应受到惯性限制，幅值增益趋于零。阻尼较小时，中间某个频率附近会出现明显的放大。

\subsection{谐振频率：使稳态振幅最大的输入频率}
\textbf{保持 $A$ 不变，使稳态输出幅值 $B$ 最大的非零输入角频率，称为谐振角频率 $\omega_r$。}因为 $B=AM$，这等价于寻找 $M(\omega)$ 的最大值。

令 $r=\omega/\omega_n$。式\eqref{eq:magnitude}的分子恒定，只需使分母根号内的量最小：
\[
  D(r)=(1-r^2)^2+4\zeta^2r^2,
  \qquad
  \frac{\dd D}{\dd r}=4r(r^2+2\zeta^2-1).
\]
令导数为零，取 $r>0$，得到 $r^2=1-2\zeta^2$。只有 $1-2\zeta^2>0$ 时才有非零实数解；结合正阻尼条件，得到
\begin{equation}\label{eq:resonance-frequency}
  \boxed{\omega_r=\omega_n\sqrt{1-2\zeta^2}},
  \qquad
  \boxed{0<\zeta<\frac{1}{\sqrt2}}.
\end{equation}
在这个范围内，$D(r)$ 先减小后增大，因此该点确实对应幅值增益的最大值。

\newpage
\subsection{谐振峰：区分峰的位置、峰的高度与出现条件}
\textbf{谐振峰是幅频曲线上在非零频率处出现的峰顶。}峰顶的横坐标是谐振频率 $\omega_r$，纵坐标是谐振峰值 $M_r$。将式\eqref{eq:resonance-frequency}代入式\eqref{eq:magnitude}，分母根号内的量化为 $4\zeta^2(1-\zeta^2)$，因此
\begin{equation}\label{eq:resonance-peak}
  \boxed{M_r=M(\omega_r)=\frac{1}{2\zeta\sqrt{1-\zeta^2}}},
  \qquad 0<\zeta<\frac{1}{\sqrt2}.
\end{equation}
谐振峰值表示输出相对输入的最大幅值比，对应的稳态输出幅值为 $B_r=AM_r$。\textbf{谐振峰不是时间响应曲线中某一时刻的波峰}，也不是阶跃响应的超调量。

\noindent
\begin{tabularx}{\linewidth}{@{}p{37mm}Y@{}}
\toprule
阻尼比 & 谐振情况 \\
\midrule
$\zeta=0$ & 理想无阻尼系统在 $\omega=\omega_n$ 的正弦激励下，振幅不断增长，没有有限振幅的稳态响应。 \\
\addlinespace[3pt]
$0<\zeta<1/\sqrt2$ & 存在有限的非零频率谐振峰；阻尼越小，峰值越高。 \\
\addlinespace[3pt]
$\zeta\geq1/\sqrt2$ & 幅值增益随频率增大而下降，没有非零频率处的谐振峰；最大增益对应零频极限。 \\
\bottomrule
\end{tabularx}
\par
因此，\textbf{欠阻尼不一定有谐振峰}：欠阻尼的范围是 $0<\zeta<1$，而有限谐振峰的条件更严格，为 $0<\zeta<1/\sqrt2\approx0.707$。

\subsection{自然频率、阻尼振荡频率与谐振频率}
\noindent
\begin{tabularx}{\linewidth}{@{}>{\raggedright\arraybackslash}p{34mm}Y>{\raggedright\arraybackslash}p{43mm}@{}}
\toprule
频率名称 & 回答的问题 & 公式与条件 \\
\midrule
无阻尼自然角频率 $\omega_n$ & 没有阻尼、没有外力时，系统自身振动多快？ & $\sqrt{k/m}$\newline $m,k>0$ \\
\addlinespace[3pt]
阻尼振荡角频率 $\omega_d$ & 欠阻尼自由响应中，衰减振荡的节奏是什么？ & $\omega_n\sqrt{1-\zeta^2}$\newline $0<\zeta<1$ \\
\addlinespace[3pt]
谐振角频率 $\omega_r$ & 持续正弦输入中，哪个输入频率使稳态振幅最大？ & $\omega_n\sqrt{1-2\zeta^2}$\newline $0<\zeta<1/\sqrt2$ \\
\bottomrule
\end{tabularx}
\par
三者的单位均为 $\mathrm{rad/s}$，换成赫兹需除以 $2\pi$。存在有限谐振峰时，$\omega_r<\omega_d<\omega_n$；阻尼很小时，它们才彼此接近。

\subsection{数值例子与适用范围}
取 $\omega_n=2\unit{rad/s}$、$\zeta=0.2$，则 $\omega_d\approx1.96\unit{rad/s}$、$\omega_r\approx1.92\unit{rad/s}$、$M_r\approx2.55$。输入频率为 $\omega_r$ 时，$A=1$ 对应 $B\approx2.55$；若 $A=2$，则 $B\approx5.10$。若改为 $\zeta=0.8$，自由响应仍有衰减振荡，但幅频曲线已无谐振峰。

\begin{keybox}
\textbf{适用范围：}本章公式针对无零点标准二阶模型。若稳态增益为正常数 $K$，则幅值增益和峰值乘以 $K$，谐振频率及其出现条件不变。自由振荡由特征根判断；谐振峰则需考察指定输入到输出的幅频响应。
\end{keybox}

\newpage
\ctexset{appendix/name={附录~,}}
\appendix
\section{特征方程与特征根}\label{app:characteristic}

本附录解释正文中的一步：为什么把 $x=e^{st}$ 代入微分方程，就能判断系统是否振荡、是否稳定？以下讨论实系数、线性、常系数的二阶系统。

\subsection{先分清三个对象}
以无外力时的小车为例：
\[
  \underbrace{m\ddot x+c\dot x+kx=0}_{\text{微分方程：未知量是函数 }x(t)}
  \quad\longrightarrow\quad
  \underbrace{p(s)=ms^2+cs+k}_{\text{特征多项式：关于 }s\text{ 的多项式}}.
\]
令 $p(s)=0$，得到\textbf{特征方程}；满足这个代数方程的数 $s_1,s_2$，称为\textbf{特征根}。根可以是实数，也可以是复数。二阶方程按重数计有两个根；“两个相同的根”称为二重根。

\subsection{为什么试取指数形式？}
指数函数的导数仍与自身成比例。取一个非零试探解 $x_h(t)=Ce^{st}$，其中 $C\ne0$、$s$ 待定，则
\[
  \dot x_h=sCe^{st},\qquad \ddot x_h=s^2Ce^{st}.
\]
代入齐次方程，得到
\[
  (ms^2+cs+k)Ce^{st}=0.
\]
由于 $Ce^{st}\ne0$，可以约去它，留下 $ms^2+cs+k=0$。因此，\textbf{特征根给出系统允许的指数变化模式}。它不是位移，而是指数中的变化率参数；若 $t$ 用秒计，$s$ 的单位为 $\mathrm{s^{-1}}$。

有输入时，先把输入项置零，得到右端为零的\textbf{齐次方程}，用它研究自由响应。对一般常系数方程
\[
  a_2\ddot y+a_1\dot y+a_0y=f(t),\quad a_2\ne0,
\]
其特征方程为 $a_2s^2+a_1s+a_0=0$。输入决定受迫响应，初始条件决定各自由响应模式的系数；它们不改变给定微分方程的特征根。

\subsection{如何求根，并写出自由响应？}
令 $\Delta=a_1^2-4a_2a_0$，由二次方程求根公式得到
\[
  \boxed{s_{1,2}=\frac{-a_1\pm\sqrt{\Delta}}{2a_2}.}
\]
\noindent
\begin{tabularx}{\linewidth}{@{}p{21mm}p{40mm}Y@{}}
\toprule
判别式 & 特征根 & 齐次通解 $y_h(t)$ \\
\midrule
$\Delta>0$ & 两个不同实根 $s_1,s_2$ & $C_1e^{s_1t}+C_2e^{s_2t}$ \\
$\Delta=0$ & 二重实根 $s_0$ & $(C_1+C_2t)e^{s_0t}$ \\
$\Delta<0$ & $\alpha\pm\mathrm{j}\beta$，$\beta>0$ & $e^{\alpha t}[C_1\cos(\beta t)+C_2\sin(\beta t)]$ \\
\bottomrule
\end{tabularx}
\par
这里 $\mathrm{j}^2=-1$，$C_1,C_2$ 由两个初始条件确定。重根时，两个相同的 $e^{s_0t}$ 不能提供两个独立解，因此第二个解要取 $te^{s_0t}$。

\newpage
\subsection{实部管增长或衰减，虚部管振荡}
把一个根写成 $s=\alpha+\mathrm{j}\beta$，利用欧拉公式：
\[
  e^{st}=e^{\alpha t}e^{\mathrm{j}\beta t}
  =e^{\alpha t}[\cos(\beta t)+\mathrm{j}\sin(\beta t)].
\]
这说明复数根并不要求物理位移是复数。实系数方程的非实根成共轭对出现，组合后就得到 A.3 节的实值正弦、余弦响应。
\begin{itemize}
  \item $\alpha<0$：该模式的包络衰减；$\alpha>0$：包络增长。
  \item $\beta\ne0$：存在振荡项，角频率为 $|\beta|$，周期为 $2\pi/|\beta|$。
  \item $\beta=0$：该指数模式没有周期性振荡；是否衰减仍由实部决定。
\end{itemize}
例如，$s=-2\pm3\mathrm{j}$ 对应
\[
  y_h(t)=e^{-2t}[C_1\cos(3t)+C_2\sin(3t)].
\]
其包络时间常数为 $1/2\unit{s}$，振荡角频率为 $3\unit{rad/s}$，周期为 $2\pi/3\approx2.09\unit{s}$。

\noindent
\begin{tabularx}{\linewidth}{@{}p{57mm}Y@{}}
\toprule
特征根在复平面中的位置 & 自由响应与稳定性 \\
\midrule
所有根的实部都小于零 & 任意初始扰动都衰减到零，称为渐近稳定；负实部重根也满足这一点。 \\
\addlinespace[3pt]
至少一个根的实部大于零 & 存在增长模式，系统不稳定。 \\
\addlinespace[3pt]
无右半平面根，但有虚轴上的根 & 需要继续检查重数。对本附录的标量二阶方程，虚轴上的单根可产生有界但不衰减的自由响应；虚轴重根可因 $t$ 因子产生无界响应。 \\
\bottomrule
\end{tabularx}
\par
\textbf{无阻尼的 $s=\pm\mathrm{j}\omega_n$ 是临界稳定的自由运动，不是渐近稳定。}同时，它并不满足“有界输入产生有界输出”：以 $\omega_n$ 持续正弦激励时会共振增长。

\subsection{放回标准二阶模型，读出阻尼性质}
对 $\omega_n>0$ 的标准二阶模型，齐次方程及特征根为
\[
  \ddot y+2\zeta\omega_n\dot y+\omega_n^2y=0,
  \qquad s_{1,2}=-\zeta\omega_n\pm\omega_n\sqrt{\zeta^2-1}.
\]
当 $0<\zeta<1$ 时，应把根号中的负数写成虚数形式：
\[
  \boxed{s_{1,2}=-\zeta\omega_n\pm\mathrm{j}\omega_n\sqrt{1-\zeta^2}
  =-\zeta\omega_n\pm\mathrm{j}\omega_d.}
\]
所以实部 $-\zeta\omega_n$ 决定衰减，虚部的绝对值 $\omega_d$ 决定自由振荡频率。$\zeta=1$ 时两根合并为 $-\omega_n$；$\zeta>1$ 时为两个负实根。\textbf{判别式判断根的类型，实部的符号判断稳定性}，两件事应分开检查。

\newpage
\subsection{完整算例：从方程求到具体运动}
取 $m=1\unit{kg}$、$c=2\unit{N\,s/m}$、$k=4\unit{N/m}$。无外力时，给定 $x(0)=x_0$、$\dot x(0)=0$：
\[
  \ddot x+2\dot x+4x=0
  \quad\Rightarrow\quad s^2+2s+4=0
  \quad\Rightarrow\quad s_{1,2}=-1\pm\mathrm{j}\sqrt3.
\]
先由根写出自由响应，再用初始条件确定系数：
\[
  x(t)=e^{-t}[C_1\cos(\sqrt3\,t)+C_2\sin(\sqrt3\,t)],
  \qquad C_1=x_0,
\]
\[
  \dot x(0)=-C_1+\sqrt3\,C_2=0
  \quad\Rightarrow\quad C_2=\frac{x_0}{\sqrt3}.
\]
因此，
\[
  \boxed{x(t)=x_0e^{-t}\left[\cos(\sqrt3\,t)+\frac{1}{\sqrt3}\sin(\sqrt3\,t)\right].}
\]
与正文参数对照：$\omega_n=2\unit{rad/s}$、$\zeta=0.5$、$\omega_d=\sqrt3\unit{rad/s}$。改变 $x_0$ 会改变振幅，但特征根仍为 $-1\pm\mathrm{j}\sqrt3$。

\subsection{它与传递函数极点、矩阵特征值有什么关系？}
\textbf{从传递函数看：}本讲义的标准模型为
\[
  G(s)=\frac{K\omega_n^2}{s^2+2\zeta\omega_ns+\omega_n^2}.
\]
当 $K\ne0$ 时，分子是非零常数，没有极零相消，分母的根就是传递函数的\textbf{极点}，也就是上述特征根。对更一般的模型，约分可能隐藏某些内部模式，不能把某个传递函数的极点不加条件地等同于全部内部特征根。

\textbf{从状态方程看：}令 $\boldsymbol z=[x,\dot x]^{\mathsf T}$，则无外力小车可以写成
\[
  \dot{\boldsymbol z}=A\boldsymbol z,\qquad
  A=\begin{bmatrix}0&1\\-k/m&-c/m\end{bmatrix}.
\]
矩阵特征值 $\lambda$ 满足
\[
  \det(\lambda I-A)=\lambda^2+\frac{c}{m}\lambda+\frac{k}{m}=0.
\]
其中 $I$ 为单位矩阵，$\det$ 表示行列式。它与小车的特征方程只差一个非零常数因子 $m$，所以两者的根相同。$s$ 和 $\lambda$ 只是不同写法下采用的符号。

\subsection{回看时记住这条推理链}
\begin{keybox}
\textbf{齐次微分方程 $\to$ 试取指数解 $\to$ 特征方程 $\to$ 特征根 $\to$ 自由响应模式。}

特征根决定允许的模式，初始条件决定各模式的系数。非实根不意味着不稳定；负实部的共轭复根对应衰减振荡。根的虚部给出自由振荡频率，\textbf{谐振频率还要从幅频响应求取}，不能直接用虚部代替。
\end{keybox}

\newpage
\section{极点图与根轨迹：从读图到调参}\label{app:poles-locus}

附录 A 把微分方程变成了特征根。这次把根画出来，再观察控制器增益改变时根的位置怎样变化。学完后，你应能解释实验页中的两幅 $s$ 平面图，并把它们与阶跃响应联系起来。

\subsection{极点图：把一个系统的根放到复平面}
对已约去公共因子的传递函数 $H(s)=N(s)/D(s)$，使分母为零的数叫\textbf{极点}，使分子为零的数叫\textbf{零点}。例如
\[
  H(s)=\frac{s+4}{s^2+4s+13}
  \quad\Rightarrow\quad
  \text{极点 }-2\pm3\mathrm{j},\qquad \text{零点 }-4.
\]
在复平面写 $s=\sigma+\mathrm{j}\omega$：横轴是实部 $\sigma$，纵轴是虚部 $\omega$；通常用 $\times$ 标极点、用 $\circ$ 标零点。所以上面的两个极点分别画在 $(-2,3)$、$(-2,-3)$，零点画在 $(-4,0)$。实系数系统的非实极点成共轭对，因此上下对称。零点会改变输出响应的形状，但不能把零点的左右位置当作判断系统稳定性的依据。

\begin{center}
\begin{tikzpicture}[x=0.72cm,y=0.72cm,>=Latex,line width=0.8pt,font=\small]
  \fill[soft] (-5,-4) rectangle (0,4.1);
  \fill[orange!6] (0,-4) rectangle (2.1,4.1);
  \draw[step=1,muted!15,very thin] (-5,-4) grid (2,4);
  \draw[->,muted] (-5.2,0)--(2.4,0) node[right,text=ink] {$\sigma$};
  \draw[->,muted] (0,-4.2)--(0,4.5) node[above,text=ink] {$\omega$};
  \node[anchor=north east,text=muted] at (-0.08,-0.07) {$0$};
  \node[anchor=south,text=accent,font=\footnotesize] at (-3.6,4.2) {左半平面：衰减};
  \node[anchor=south,text=muted,font=\footnotesize] at (1.1,4.2) {增长};
  \draw[dashed,muted] (-2,0)--(-2,3)--(0,3);
  \draw[dashed,muted] (-2,0)--(-2,-3)--(0,-3);
  \draw[accent,thick] (0,0)--(-2,3);
  \node[rotate=-56,fill=soft,inner sep=1.5pt,text=accent] at (-1.15,1.75) {$\omega_n=\sqrt{13}$};
  \node[accent,font=\Large] at (-2,3) {$\times$};
  \node[accent,font=\Large] at (-2,-3) {$\times$};
  \node[anchor=south east,fill=soft,inner sep=2pt,text=ink] at (-2.15,3.1) {$-2+3\mathrm{j}$};
  \node[anchor=north east,fill=soft,inner sep=2pt,text=ink] at (-2.15,-3.1) {$-2-3\mathrm{j}$};
  \draw[accent,thick,fill=white] (-4,0) circle (0.10);
  \node[below,text=ink] at (-4,-0.15) {$-4$};
  \node[below,text=ink] at (-2,-0.15) {$-2$};
  \node[right,text=ink] at (0.05,3) {$3$};
  \node[right,text=ink] at (0.05,-3) {$-3$};
  \node[anchor=north west,text width=5.4cm,align=left,text=ink] at (3.5,3.4) {
    \textbf{先读图例，再读坐标}\\[5pt]
    $\times$：极点 $-2\pm3\mathrm{j}$\\[4pt]
    $\circ$：零点 $-4$\\[8pt]
    横坐标 $-2$：包络按 $e^{-2t}$ 衰减。\\[5pt]
    纵坐标 $\pm3$：振荡角频率为 $3\,\mathrm{rad/s}$。\\[5pt]
    两轴等比例；坐标单位为 $\mathrm{s^{-1}}$。横轴不是时间。
  };
\end{tikzpicture}
\end{center}


\textbf{这两条坐标轴都不是时间轴。}图上的一个点表示一个复数变化率，不是某时刻的位置、速度或输出大小；两点也不是小车往返运动的两个端点。

\subsection{读一个点，就读出一种变化模式}
沿用附录 A，$s=-2\pm3\mathrm{j}$ 对应自由响应中的
\[
  e^{-2t}\big[C_1\cos(3t)+C_2\sin(3t)\big].
\]
向左的距离 $2$ 给出包络 $e^{-2t}$，时间常数为 $\tau=1/2\unit{s}$；虚部大小 $3$ 给出振荡角频率 $3\unit{rad/s}$，周期为 $2\pi/3\approx2.09\unit{s}$。对于这一对二阶极点，还可读出
\[
  \omega_n=\sqrt{\sigma^2+\omega^2}=\sqrt{13}\unit{rad/s},
  \qquad \zeta=\frac{-\sigma}{\omega_n}=\frac{2}{\sqrt{13}}\approx0.555.
\]
注意 $\omega_n$ 是点到原点的距离，$|\omega|$ 才是自由振荡角频率；两者通常不相等。

\newpage
\subsection{先看稳定性，再区分两种图}
以下讨论有限维、连续时间线性时不变系统。全部内部特征根都在左半平面，初始扰动最终衰减；有根进入右半平面，就存在增长模式。虚轴是边界：虚轴上的单根可产生不衰减的自由运动，并不能据此称系统渐近稳定；虚轴重根还要检查是否产生随时间增长的项。未被相消的虚轴极点也不满足“有界输入产生有界输出”，例如持续同频正弦输入可导致共振增长。

\textbf{不要用极零相消掩盖内部不稳定。}最简传递函数只反映该输入到输出的关系，约分可能隐藏内部模式。本实验的“当前极点图”保留闭环特征多项式的全部根；发生精确相消时，它与约分后的传递函数极点图可能不同。

\begin{keybox}
\textbf{极点图回答：这组参数下，根在哪里？}

\textbf{根轨迹回答：固定开环形状，只连续改变一个整体增益，闭环根经过哪些位置？}
\end{keybox}

极点图是一张固定参数的快照；根轨迹把同一族系统的许多快照连起来。轨迹上的箭头表示增益增大的方向，\textbf{不是时间流逝的方向}。选定一个增益后，实际系统的极点就固定下来，不会沿着图中的曲线随时间移动。

\subsection{根轨迹为什么求的是闭环特征方程？}
设固定的开环形状为 $L_0(s)=N(s)/D(s)$，在其前面乘整体增益 $\kappa\ge0$。单位负反馈给出
\[
  T_\kappa(s)=\frac{\kappa L_0(s)}{1+\kappa L_0(s)}
  =\frac{\kappa N(s)}{D(s)+\kappa N(s)}.
\]
因此，每选一个 $\kappa$，就求一次\textbf{闭环特征方程}
\[
  \boxed{D(s)+\kappa N(s)=0.}
\]
根轨迹描绘这些根随 $\kappa$ 的变化。开环极点是 $D(s)=0$ 的根，闭环根则随反馈增益变化，不能把二者混为一谈。

在 $N,D$ 无公共因子、$L_0$ 严格真有理的通常情形，若分母次数为 $n$、分子次数为 $m<n$，根轨迹有 $n$ 条分支（按重数计）。当 $\kappa\to0$，它们趋近开环极点；当 $\kappa\to\infty$，其中 $m$ 条趋近有限开环零点，其余走向无穷远。$\kappa=0$ 时 $T_0(s)=0$，所以起点应理解为\textbf{特征根的极限}，而不是零传递函数中可见的输出极点。

\newpage
\subsection{亲手算一条根轨迹}
先选比 PID 更容易手算的对象 $G_0(s)=1/[(s+1)(s+3)]$，控制器只有比例增益 $C(s)=\kappa$，采用单位负反馈。于是
\[
  T_\kappa(s)=\frac{\kappa}{s^2+4s+3+\kappa},
  \qquad s_{1,2}=-2\pm\sqrt{1-\kappa}.
\]
只需替换五个增益，就能自己画出轨迹的骨架：

\noindent
\begin{tabularx}{\linewidth}{@{}p{22mm}p{46mm}Y@{}}
\toprule
$\kappa$ & 闭环特征根 & 读图结论 \\
\midrule
$0$（极限） & $-1,\,-3$ & 从两个开环极点出发 \\
$0.75$ & $-1.5,\,-2.5$ & 沿实轴相向靠近 \\
$1$ & $-2$（二重根） & 在实轴上汇合 \\
$2$ & $-2\pm\mathrm{j}$ & 离开实轴，上下对称 \\
$5$ & $-2\pm2\mathrm{j}$ & 虚部增大，实部仍为 $-2$ \\
\bottomrule
\end{tabularx}
\par

\begin{center}
\begin{tikzpicture}[x=1cm,y=1cm,>=Latex,line width=0.8pt,font=\small]
  \fill[soft] (-4.3,-2.9) rectangle (0,2.9);
  \draw[step=1,muted!15,very thin] (-4,-2) grid (0,2);
  \draw[->,muted] (-4.4,0)--(0.7,0) node[right,text=ink] {$\sigma$};
  \draw[->,muted] (0,-3)--(0,3.2) node[above,text=ink] {$\omega$};
  \node[below right,text=ink] at (0,0) {$0$};
  \draw[accent,very thick] (-3,0)--(-1,0);
  \draw[->,accent,very thick] (-2,0)--(-2,2.8);
  \draw[->,accent,very thick] (-2,0)--(-2,-2.8);
  \draw[->,accent,very thick] (-2.9,0)--(-2.45,0);
  \draw[->,accent,very thick] (-1.1,0)--(-1.55,0);
  \node[text=ink,font=\Large] at (-3,0) {$\times$};
  \node[text=ink,font=\Large] at (-1,0) {$\times$};
  \node[below=5pt,text=ink] at (-3,0) {$-3$};
  \node[below=5pt,text=ink] at (-1,0) {$-1$};
  \fill[accent] (-2,0) circle (0.045);
  \node[anchor=south east,text=ink,fill=soft,inner sep=2pt] at (-2.18,0.18) {$\kappa=1$};
  \node[anchor=north east,text=muted] at (-2.1,-0.15) {$-2$};
  \fill[accent] (-2,1) circle (0.065);
  \fill[accent] (-2,-1) circle (0.065);
  \node[right,text=ink] at (-1.88,1) {$\kappa=2$};
  \node[right,text=ink] at (-1.88,-1) {$\kappa=2$};
  \fill[accent] (-2,2) circle (0.065);
  \fill[accent] (-2,-2) circle (0.065);
  \node[right,text=ink] at (-1.88,2) {$\kappa=5$};
  \node[right,text=ink] at (-1.88,-2) {$\kappa=5$};
  \node[anchor=north west,text width=5.4cm,align=left,text=ink] at (1.5,2.65) {
    \textbf{沿箭头增大 $\kappa$}\\[6pt]
    1. 从开环极点 $-3$、$-1$ 出发。\\[6pt]
    2. 在 $-2$ 相遇，形成二重根。\\[6pt]
    3. 变成共轭复根，沿竖直线离开实轴。\\[8pt]
    $\times$ 是开环起点；蓝线上的点是不同增益下的闭环根。\\[6pt]
    图只画有限区域，向上下的分支继续延伸。
  };
\end{tikzpicture}
\end{center}


在 $0<\kappa<1$ 时，靠近虚轴的慢极点向左移动；越过 $\kappa=1$ 后，两根沿 $\sigma=-2$ 的竖线分开。此时
\[
  \tau=\frac12\unit{s},\qquad
  \omega_d=\sqrt{\kappa-1},\qquad
  \zeta=\frac{2}{\sqrt{3+\kappa}}.
\]
因此，继续增加增益会使振荡频率升高、阻尼比下降，包络时间常数却不再缩短。增益越大并不保证收敛越快；该系统的单位阶跃终值 $\kappa/(3+\kappa)$ 也随增益变化。判断跟踪效果时，还要一起检查稳态误差和完整响应。

\newpage
\subsection{把读图方法放回现有实验}
实验对象是 $G(s)=4/(s^2+1.2s+4)$。仅用 P 控制时，令 $K_i=K_d=0$，有
\[
  T(s)=\frac{4K_p}{s^2+1.2s+4+4K_p},
  \qquad s_{1,2}=-0.6\pm\mathrm{j}\sqrt{3.64+4K_p}
  \quad(K_p\ge0).
\]
当 $K_p=1,3,8$ 时，极点分别约为 $-0.6\pm2.764\mathrm{j}$、$-0.6\pm3.955\mathrm{j}$、$-0.6\pm5.970\mathrm{j}$。它们只是上下分开，没有向左移动；振荡包络的时间常数始终为 $1/0.6\approx1.67\unit{s}$，单位阶跃终值则依次为 $0.5,0.75,8/9$。

\textbf{实验页里的 $\kappa$ 与 $K_p$ 不是两个名称相同的滑块。}当前 PID 为
\[
  C(s)=K_p+\frac{K_i}{s}+K_ds,
\]
根轨迹固定当前的 $K_p:K_i:K_d$，扫描的是 $\kappa C(s)G(s)$，即三个增益一起乘 $\kappa$。图中实心点标记 $\kappa=1$，正好对应当前控制器，应与“当前极点图”的特征根重合。

单独改变 $K_p$、$K_i$ 或 $K_d$，通常会改变控制器零点和开环形状，所以软件会重新计算一族根轨迹。\textbf{纯 P 是一个例外：}正的 $K_p$ 只改变整体比例，此时有效增益是 $\kappa K_p$；若都扫描到无穷大，轨迹的几何形状相同，只是增益标尺不同。实验只画所选的有限 $\kappa$ 范围，所以显示出的轨迹长度仍会改变。

\begin{keybox}
极点告诉你可能出现哪些衰减或振荡模式，\textbf{不能单独给出完整阶跃响应}。零点、各模式的系数（残差）、输入和初始条件都会影响输出；PID 还可能引入第三个极点。因此，配置一对目标极点后，不能直接把标准二阶超调公式当作完整 PID 闭环的仿真结果。
\end{keybox}

\newpage
\subsection{打开实验页，完成三步观察}
进入\href{https://youngx-ray.github.io/labs/second-order-pid/}{PID 极点配置实验}。每次先预测极点移动方向，再改参数，最后核对阶跃曲线。可以点击“固定当前响应作对照”，保留改动前的曲线。

\textbf{第一步：只加比例，检查“频率变快”和“衰减变快”。}
点击“P”，确认 $K_i=K_d=0$，依次把“比例 $K_p$”数值设为 $1,3,8$。记录极点实部、虚部与阶跃终值。你应看到实部始终为 $-0.6$，波峰间隔缩短，但包络没有因实部左移而加快衰减。根轨迹的实心点始终表示当前 $\kappa=1$；把右上角 $\kappa$ 上限从 $1$ 改为 $5$，只是多看一段轨迹，不会改动当前阶跃响应。

\textbf{第二步：加入微分，让实部发生变化。}
点击“PD”，再设 $K_p=2$、$K_i=0$，依次尝试 $K_d=0,0.3,0.75$。此时闭环特征方程为
\[
  s^2+(1.2+4K_d)s+12=0.
\]
这三组参数下，两根仍为共轭复根，实部依次是 $-0.6,-1.2,-2.1$，对应越来越快的包络衰减。观察零点也随之变化，再查看实际超调。继续无限增大微分并不保证更快：进入过阻尼后，可能出现一个靠近原点的慢极点。

\textbf{第三步：加入积分，亲眼看极点穿过虚轴。}
点击“PI”，设 $K_p=2$、$K_d=0$，把 $K_i$ 依次设为 $3,3.6,4$。特征方程及稳定条件为
\[
  s^3+1.2s^2+12s+4K_i=0,
  \qquad 0<K_i<3.6.
\]
预期结果：$K_i=3$ 时三根都在左半平面，阶跃输出最终趋于 $1$；$K_i=3.6$ 时分母为 $(s+1.2)(s^2+12)$，有一对虚轴根 $\pm\mathrm{j}\sqrt{12}$，振荡不再衰减；$K_i=4$ 时一对根进入右半平面，出现增长振荡。把“观察时长”改为 $60\unit{s}$ 更容易看出区别。\textbf{稳态误差为零的结论以闭环稳定为前提}，不能套用于后两种情况。

\newpage
\subsection{四道自测：先判断，再核对}
\begin{itemize}
  \item \textbf{题 1：}图上极点为 $-1\pm2\mathrm{j}$，零点为 $-5$。分别画在哪里？振荡包络时间常数和周期是多少？
  \item \textbf{题 2：}两组极点分别为 $-1\pm2\mathrm{j}$ 与 $-1\pm5\mathrm{j}$。哪组自由振荡更快？哪组指数包络衰减更快？
  \item \textbf{题 3：}在 B.5 的手算系统中，取 $\kappa=10$，求根并判断稳定性。轨迹上的点会随时间跑向上方吗？
  \item \textbf{题 4：}实验中固定 $K_p=2,K_d=0$，把 $K_i$ 从 $3$ 改为 $4$。这是否只是沿原先固定 PID 比例的根轨迹移动？新闭环是否稳定？
\end{itemize}

\textbf{答案 1：}在 $(-1,2)$ 与 $(-1,-2)$ 画 $\times$，在 $(-5,0)$ 画 $\circ$。极点对应包络 $e^{-t}$，所以 $\tau=1\unit{s}$，周期为 $2\pi/2=\pi\unit{s}$。零点不应被当作另一个自由响应模式。

\textbf{答案 2：}第二组振荡更快，周期从 $\pi$ 缩短到 $2\pi/5$；两组指数包络都是 $e^{-t}$，衰减时间常数相同。但完整输出的振幅、超调和实际进入调节带的时间，还取决于传递函数及输入。

\textbf{答案 3：}$s=-2\pm\sqrt{1-10}=-2\pm3\mathrm{j}$，两根都在左半平面，闭环稳定。选定 $\kappa=10$ 后极点固定；响应随时间变化，极点坐标不会沿根轨迹运动。

\textbf{答案 4：}不是。只改 $K_i$ 改变了 PID 比例和控制器零点，需重新生成根轨迹。原条件要求 $K_i<3.6$，因此 $K_i=4$ 的闭环不稳定。应同时看到当前极点进入右半平面、阶跃响应振荡增长。

\begin{keybox}
\textbf{读图顺序：先认对象和坐标，再看左右与上下。}

确认图画的是开环极点、闭环特征根，还是闭环零点；横看衰减或增长，纵看振荡频率。若有轨迹，再问“固定了什么、扫描哪个增益、当前点在哪里”。最后回到阶跃响应，核对稳态误差、超调和收敛过程。
\end{keybox}

符号约定与根轨迹定义可参见 MathWorks 官方文档：\href{https://www.mathworks.com/help/control/ref/lti.pzmap.html}{极点与零点图 pzmap}、\href{https://www.mathworks.com/help/control/ref/dynamicsystem.rlocus.html}{根轨迹 rlocus}。

\end{document}
